\documentclass[english, parskip=full]{scrartcl}
\usepackage[english]{babel}  % Sprache auf Deutsch 
\usepackage[utf8]{inputenc}  % Kodierung der Datei
\usepackage[left=2cm, right=2cm, top=2cm, bottom=3cm, footskip=1.5cm]{geometry}
\input{myscrarticle}

\newcommand*{\titel}{Matsubara Summation}
\newcommand*{\autor}{Joachim Schwardt}

\hypersetup{pdfauthor={\autor}, pdftitle={\titel}}

%\titlehead{titlehead}
\subject{Mathematical Methods}
\title{\titel}
\author{\autor}
\date{\today}

\begin{document}

%\begin{titlepage}
\maketitle  % Titel setzen
%\tableofcontents  % Inhaltsverzeichnis setzen
%\end{titlepage}


\section{Formalism}
Let $g(z)$ be a complex function satisfying the asymptotic relation $|g(z)| \sim \frac{1}{|z|^{1+\epsilon}}$ for $\epsilon > 0$ and let $z_k$ for $k=1,\dots,K$ denote the $K$ poles of $g$.
Then for $f_\text{F}(z) = \frac{\powe{\beta z} - 1}{\powe{\beta z} + 1}$ and $f_\text{B}(z) = \frac{\powe{\beta z} + 1}{\powe{\beta z} - 1}$ we have the identities
\begin{align}
0 &= \frac{1}{2\pi\i} \oint_{|z| \rightarrow\infty} g(z) f_{\text{F,B}} \,\mathrm{d}z = \frac{2}{\beta} \sum_{n\in\mathbb{Z}} g(\i\omega_n^\Text{F,B}) + \frac{1}{2\pi\i}\sum_{k=1}^K \Res\klb*{g\circ f_\text{F,B} }(z_k)
\end{align}
where $\omega_n^\text{F} = \frac{(2n+1)\pi}{\beta}$ and $\omega_n^\text{B} = \frac{2\pi n}{\beta}$ are the fermionic and bosonic Matsubara frequencies.
\section{List of problems}
\subsection{Basel problem}
We want to solve the generalized Basel problem
\begin{align}
S_{N} &= \sum_{n\ge 1} \frac{1}{n^{N}}.
\end{align}
We can not immediately translate this into a Matsubara sum, since the $n=0$-term diverges.
\subsubsection{Higher-order pole}
Introducing a regularization $\epsilon > 0$ we may define
\begin{align}
S_N(\epsilon) &= \frac{-1}{2(-\epsilon)^N} + \frac{1}{2}\sum_{n\in\mathbb{Z}} \frac{1}{(n - \epsilon)^N}
\end{align}
where $S_N(0) = S_N$ for even $N$.
We will choose the bosonic frequencies for $\beta = 2$ and define
\begin{align}
g(z) &= \frac{1}{(z - \epsilon)^N}.
\end{align}
This function has an $N$-th order pole at $z=\epsilon$ and thus
\begin{align}
-\sum_{n\in\mathbb{Z}} g(\i\omega_n^\text{B}) &= \frac{1}{(N-1)!} \lim_{z\rightarrow \epsilon} \partial_z^{N-1} (z-\epsilon)^N g(z) f_\text{B}(z) = \frac{1}{(N-1)!} \lim_{z\rightarrow \epsilon} \partial_z^{N-1} f_\text{B}(z).
\end{align}
Again assuming even $N$ and as $\epsilon\rightarrow 0$ the left hand side simplifies to
\begin{align}
- \frac{1}{(-\epsilon)^N} - 2 \sum_{n\ge 1} \frac{1}{\kla*{\pi\i n}^N} = -\frac{1}{(-\epsilon)^N} - 2\kla*{\frac{1}{\pi\i}}^N S_N.
\end{align}
Collecting terms gives
\begin{align}
S_N &= -\lim_{\epsilon \rightarrow 0} \kla*{\pi\i}^N \frac{1}{2} \klb*{ \frac{1}{(-\epsilon)^N} +  \frac{1}{(N-1)!} f_\text{B}^{(N-1)}(\epsilon)  }.
\end{align}
Using $f_\text{B}(z) = 1 + \frac{2}{\powe{2 z} - 1} = \coth(z)$ simplifies the derivatives a bit.
However, applying this formula for general $N$ is still somewhat involved.
Note that we can rewrite the expression as
\begin{align}
S_N &= -\frac{(\pi\i)^N}{2(N-1)!} \lim_{\epsilon \rightarrow 0} \partial_\epsilon^{N-1} \klb*{ \coth(\epsilon) - \frac{1}{\epsilon} }.
\end{align}
In this form it becomes obvious that given a series representation $\coth(z) = \frac{1}{z} + \sum_{n\ge 1} a_n z^{2n-1}$ for small $z$ results in
\begin{align}
S_N &= -\frac{(\pi\i)^N}{2(N-1)!} \lim_{\epsilon \rightarrow 0} \partial_\epsilon^{N-1} \sum_{n\ge 0} a_n \epsilon^{2n-1} = \pi^N \frac{-(-1)^\frac{N}{2} a_{\frac{N}{2}}}{2(N-1)!}.
\end{align}
According to literature the coefficients are $a_n = \frac{4^n B_{2n}}{(2n)!}$.\\
For the simplest case of $N=2$ we find
\begin{align}
S_N &= -\lim_{\epsilon \rightarrow 0} \kla*{\pi\i}^N \frac{1}{2} \klb*{ \frac{1}{\epsilon^2} -\frac{1}{\sinh^2 (\epsilon)} } = \frac{\pi^2}{2} \lim_{\epsilon\rightarrow 0} \klb*{\frac{1}{\epsilon^2} - \frac{1}{\epsilon^2 + \frac{\epsilon^4}{3}}} = \frac{\pi^2}{2} \lim_{\epsilon\rightarrow 0} \frac{\frac{\epsilon^4}{3}}{\epsilon^2 \kla*{\epsilon^2 + \frac{\epsilon^4}{3}} } = \frac{\pi^2}{6}.
\end{align}
It is interesting how this approach fails for odd $N$. 
If one replaces $z$ with $|z|$ in the definition of $g(z)$, the positive and negative indices would no longer cancel each other out, and one would indeed get a result.
The problem is, however, that the resulting complex function does not have a simple pole at $z=\epsilon$ but rather a continuum of poles with $|z|=\epsilon$, i.e. it is not analytic.
Solving this problem would require a different contour, or somehow including another scaling factor to suppress the negative Matsubara frequencies.
\subsubsection{Multiple simple poles}
Another way to introduce a regularization is to define
\begin{align}
S_N(\epsilon) &= \frac{1}{2\epsilon^N} + \frac{1}{2} \sum_{n\in\mathbb{Z}} \frac{1}{n^N - \epsilon^N}.
\end{align}
Instead of a single pole of $N$-th order we will instead have to treat $N$ simple poles.
Inspired by this formulation we define
\begin{align}
g(z) &= \frac{1}{z^N - \epsilon^N}
\end{align}
which has poles at $z_k=\epsilon \powe{2\pi\i \frac{k}{N}}$ for $k=0,\dots,N-1$.
For $\beta = 2$ we find
\begin{align}
-\sum_{n\in\mathbb{Z}} g(\i\omega_n^\text{B}) &= \sum_{k=0}^{N-1} \lim_{z\rightarrow z_k} f_\text{B}(z) \frac{z - z_k}{z^N - \epsilon^N} = \sum_{k=0}^{N-1} \coth (z_k) \frac{1}{N z_k^{N-1}}.
\end{align}
The left hand side gives
\begin{align}
\frac{1}{\epsilon^N} - 2\sum_{n\ge 1} \frac{1}{(\pi\i n)^N - \epsilon^N} &\sim \frac{1}{\epsilon^N} - \frac{2}{(\pi\i)^N} S_N
\end{align}
as $\epsilon\rightarrow 0$.
Rearranging then yields
\begin{align}
S_N &= \frac{(\pi\i)^N}{2} \lim_{\epsilon \rightarrow 0} \klb*{\frac{1}{\epsilon^N} - \frac{1}{N} \sum_{k=0}^{N-1} \frac{\coth(z_k)}{z_k^{N-1}} }.
\end{align}
More explicitly the second term becomes
\begin{align}
\frac{1}{N} \sum_{k=0}^{N-1} \frac{\coth \kla*{\epsilon \powe{2\pi\i \frac{k}{N}}} }{\epsilon^{N-1} \powe{-2\pi\i \frac{k}{N}} } &= \frac{1}{N\epsilon^N} \sum_{k=0}^{N-1} \epsilon \powe{2\pi\i \frac{k}{N}} \coth \kla*{\epsilon \powe{2\pi\i \frac{k}{N}}}.
\end{align}
Working out this expression is even more involved than before, and ultimately again results in computing the $N$-th derivative of $\coth$ at zero.
\subsection{Asymptotically discontinuous functions}
\subsubsection{Step function}
Let us consider again the Basel problem
\begin{align}
S_2 &= \sum_{n\ge 1} \frac{1}{n^2}.
\end{align}
One obvious way of formally rewriting this as a sum over all integers is to introduce a step function
\begin{align}
S_2 &= \sum_{n\in\mathbb{Z}} \Theta(n - a) \frac{1}{n^2}
\end{align}
where $a\in (0,1)$.
Unfortunately the Matsubara trick does require an analytic function, which the step function is obviously not.
However, we can for example write it as the limit
\begin{align}
\Theta(z) &= \lim_{\epsilon\rightarrow 0} \Theta_\epsilon(z) = \lim_{\epsilon\rightarrow 0} \klb*{\frac{1}{2} + \frac{1}{\pi} \arctan \kla*{\frac{z}{\epsilon}}}.
\end{align}
This representation is useful since it does not introduce any additional poles.
Thus the parametrized version of our series corresponds to 
\begin{align}
g(z) &= \frac{\Theta_\epsilon\kla*{\frac{-\i z}{2\pi}} }{z^2}.
\end{align}
Note that this function does have a second-order pole at $z=0$.
Applying the fermionic Matsubara summation for $\beta=2$ gives
\begin{align}
-\sum_{n\in\mathbb{Z}} g(\i\omega_n^\text{F}) &= \sum_{n\in\mathbb{Z}} \frac{4}{\pi^2} \frac{\Theta_\epsilon \kla*{n + \frac{1}{2}} }{(2n+1)^2} =  \lim_{z\rightarrow 0} \partial_z \tanh(z) \Theta_\epsilon\kla*{\frac{-\i z}{2\pi}}.
\end{align}
We can safely take the limit on the left hand side 
\begin{align}
\frac{4}{\pi^2}\sum_{n\ge 1} \frac{1}{(2n - 1)^2} &= \lim_{\epsilon \rightarrow 0} \lim_{z\rightarrow 0} \partial_z \klb*{\frac{1}{2} \tanh(z) + \frac{1}{\pi} \tanh(z) \arctan\kla*{\frac{-\i z}{2\pi\epsilon}}  }.
\end{align}
The term in brackets is analytic near the origin, and since we only take the first derivative, it is sufficient to compute the Taylor expansion to linear order.
This gives
\begin{align}
\frac{4}{\pi^2} \sum_{n\ge 1} \frac{1}{(2n-1)^2} &= \lim_{\epsilon \rightarrow 0} \lim_{z\rightarrow 0} \partial_z \klb*{\frac{1}{2} z + \frac{1}{\pi} z \frac{-\i z}{2\pi\epsilon} } = \frac{1}{2}
\end{align}
and thus $\sum_{n\ge 1} \frac{1}{(2n-1)^2} = \frac{\pi^2}{8}$.
We can rewrite the general Basel problem using
\begin{align}
\sum_{n\ge 1} \frac{1}{n^N} &= \sum_{n\ge 1} \frac{1}{(2n-1)^N} + \sum_{n\ge 1} \frac{1}{(2n)^N} = \sum_{n\ge 1} \frac{1}{(2n-1)^N} + \frac{1}{2^N} \sum_{n\ge 1} \frac{1}{n^N}.
\end{align}
Solving for the term of interest gives $\sum_{n\ge 1} \frac{1}{n^N} = \frac{2^N}{2^N - 1} \sum_{n\ge 1} \frac{1}{(2n-1)^N} $.
In our case $N=2$ and therefore $\sum_{n\ge 1} \frac{1}{n^2} = \frac{4}{3} \sum_{n\ge 1} \frac{1}{(2n-1)^2} = \frac{4}{3} \frac{\pi^2}{8} = \frac{\pi^2}{6} $.\\
The interesting point to note here is that we did not have to make the explicit restriction to even $N$.
In fact, for $N=3$ the same procedure for a analogous $g(z)$ gives
\begin{align}
-\sum_{n\in\mathbb{Z}} \kla*{\frac{2}{\pi\i}}^3 \frac{\Theta_\epsilon\kla*{n + \frac{1}{2}} }{(2n+1)^3} &= \frac{1}{2} \lim_{z\rightarrow 0} \partial_z^2 \tanh(z) \Theta_\epsilon\kla*{\frac{-\i z}{2\pi}}.
\end{align}
Rearranging yields
\begin{align}
-\frac{8\i}{\pi^3} \sum_{n\ge 1} \frac{1}{(2n-1)^3} &= \frac{1}{2} \lim_{\epsilon \rightarrow 0} \lim_{z\rightarrow 0} \partial_z^2 \klb*{\frac{1}{2} \tanh(z) + \frac{1}{\pi} \tanh(z) \arctan\kla*{\frac{-\i z}{2\pi\epsilon}} }.
\end{align}
This time we have to expand the term to second order, which coincidentally is completely equivalent to the previous expansion
\begin{align}
-\frac{8\i}{\pi^3} \sum_{n\ge 1} \frac{1}{(2n-1)^3} &= \frac{1}{2} \lim_{\epsilon \rightarrow 0} \lim_{z\rightarrow 0} \partial_z^2 \klb*{\frac{1}{2} z + \frac{1}{\pi} z \arctan\kla*{\frac{-\i z}{2\pi\epsilon}} } = \frac{1}{2\pi} \lim_{\epsilon \rightarrow 0} \frac{2 \frac{-\i}{2\pi\epsilon}}{\kla*{\frac{-\i 0^+}{2\pi\epsilon}}^2 + 1} = ???
\end{align}
NEW ATTEMPT:\\
Consider 
\begin{align}
\Theta(z) &= \lim_{\epsilon \rightarrow 0} \Theta_\epsilon(z) = \lim_{\epsilon \rightarrow 0} \frac{1}{2} \klb*{1 + \frac{z}{\sqrt{z^2 + \epsilon^2}}}.
\end{align}
Then we have
\begin{align}
-\sum_{n\in\mathbb{Z}} \kla*{\frac{2}{\pi\i}}^3 \frac{\Theta_\epsilon\kla*{n + \frac{1}{2}} }{(2n+1)^3} &= ...
\end{align}
also not analytic, same problem as before...\\
END ATTEMPT.


%\begin{thebibliography}{9}
%\bibitem{example} description, \url{https://www.exmapleurl}, day.\,month~2021.
%\end{thebibliography}

\end{document}